\SourceLeaf{322}
comparare, quam per $X''$ exprimemus. Comparando valores coëfficientium
$\gamma,\gamma',\delta',\gamma'',\delta''$ etc. in artt. 20, 22 traditos cum valoribus coëfficientium
$\varepsilon,\varepsilon',\varepsilon''$ etc. modo inventis, invenitur
\[
\begin{aligned}
&\varepsilon=\gamma\\[5pt]
&\varepsilon'=\gamma'-\delta'\operatorname{cotang}\mu a,
\\[5pt]
&\varepsilon^{\mu-1}=\delta'\operatorname{cosec}\mu a\\[5pt]
&\varepsilon''=\gamma''-\delta''\operatorname{cotang}\mu a,
\\[5pt]
&\varepsilon^{\mu-2}=\delta''\operatorname{cosec}\mu a\\[5pt]
&\varepsilon'''=\gamma'''-\delta'''\operatorname{cotang}\mu a,
\\[5pt]
&\varepsilon^{\mu-3}=\delta'''\operatorname{cosec}\mu a\\[5pt]
&\text{etc. usque ad}\\[5pt]
&\left(\tfrac12\right)\varepsilon^m=\gamma^m-\delta^m\operatorname{cotang}\mu a,
\\[5pt]
&\left(\tfrac12\right)\varepsilon^{\mu-m}=\delta^m\operatorname{cosec}\mu a
\end{aligned}
\]
Factor $\tfrac12$ ipsis $\varepsilon^m$ et $\varepsilon^{\mu-m}$ praepositus tunc tantummodo valet, quando $\mu$ par
est, omittique debet, quando $\mu$ est impar. Quo pacto, in casu posteriore, ubi
$\mu-m=m+1$ manifesto omnes $\mu$ coëfficientes $\varepsilon,\varepsilon',\varepsilon''$ etc. per $\gamma,\gamma',\delta'$ etc.
determinati sunt; in priore vero, ubi $m=\mu-m$, pro coëfficiente $\varepsilon^{\frac12\mu}$ valores
duos $2\gamma^{\frac12\mu}-2\delta^{\frac12\mu}\cot\mu a$, $2\delta^{\frac12\mu}\operatorname{cosec}\mu a$ habemus, quorum aequalitatem ex aequatione
$\gamma^{\frac12\mu}\sin\tfrac12\mu a=\delta^{\frac12\mu}\cos\tfrac12\mu a$ facile perspicere, et pro quibus igitur etiam
hunc $\gamma^{\frac12\mu}-\delta^{\frac12\mu}\operatorname{cotg}\mu a+\delta^{\frac12\mu}\operatorname{cosec}\mu a$ adoptare possumus.

Hinc facile deducitur, esse in utroque casu
\[
\begin{aligned}
&X''-(X'-\delta'\sin x-\delta''\sin2x-\delta^m\sin3x-\text{etc.}-\delta^m\sin mx)\\[5pt]
&=-\operatorname{cotg}\mu a(\delta'\cos x+\delta''\cos2x+\delta'''\cos3x+\mathinner{.\,.\,.}+\delta^m\cos mx)\\[5pt]
&\quad+\operatorname{cosec}\mu a(\delta'\cos(\mu-1)x+\delta''\cos(\mu-2)x+\mathinner{.\,.\,.}+\delta^m\cos(\mu-m)x)
\end{aligned}
\]
Hinc patet, $X''$ ex $X'$ deduci, si pro quovis termino $\delta^\lambda\sin\lambda x$ substituatur
\[
\begin{aligned}
&-\operatorname{cotg}\mu a\,\delta^\lambda\cos\lambda x+\operatorname{cosec}\mu a\,\delta^\lambda\cos(\mu-\lambda)x
\end{aligned}
\]
sive
\[
\begin{aligned}
&\frac{-\delta^\lambda\cos\mu a\cos\lambda x+\delta^\lambda\cos(\mu-\lambda)x}{\sin\mu a}
\end{aligned}
\]
cuius praecepti comparationem cum iis, quae in art. 24 explicata sunt, lectoribus
linquimus.

\InterpolationArticle{32}
Si functio $X$ cum $\cos nx$ non abrumpitur, sed ulterius excurrit, terminis
sequentibus per $\alpha^\mu\cos\mu x+\alpha^{\mu+1}\cos(\mu+1)x+\text{etc.}$ expressis, habebimus
\SourceLeaf{323}
\[
\begin{aligned}
&\varepsilon=\alpha+\alpha^\mu\cos\mu a+\alpha^{2\mu}\cos2\mu a+\text{etc.}\\[5pt]
&\varepsilon'=\frac{1}{\sin\mu a}\{\alpha'\sin\mu a+\alpha^{\mu+1}\sin2\mu a+\alpha^{2\mu+1}\sin3\mu a+\text{etc.}\\[5pt]
&\qquad-\alpha^{2\mu-1}\sin\mu a-\alpha^{3\mu-1}\sin2\mu a-\alpha^{4\mu-1}\sin3\mu a-\text{etc.}\}\\[5pt]
&\varepsilon''=\frac{1}{\sin\mu a}\{\alpha''\sin\mu a+\alpha^{\mu+2}\sin2\mu a+\alpha^{2\mu+2}\sin3\mu+\text{etc.}\\[5pt]
&\qquad-\alpha^{2\mu-2}\sin\mu a-\alpha^{3\mu-2}\sin2\mu a-\alpha^{4\mu-2}\sin3\mu a-\text{etc.}\}\\[5pt]
&\text{etc.}\\[5pt]
&\varepsilon^{\rho-1}=\frac{1}{\sin\mu a}\{\alpha^{\mu-1}\sin\mu a+\alpha^{2\mu-1}\sin2\mu a+\alpha^{3\mu-1}\sin3\mu a+\text{etc.}\\[5pt]
&\qquad-\alpha^{\mu+1}\sin\mu a-\alpha^{2\mu+1}\sin2\mu a-\alpha^{3\mu+1}\sin3\mu a-\text{etc.}\}
\end{aligned}
\]
ex quibus formulis iudicari potest, quatenus differentiam coëfficientium $\varepsilon,\varepsilon',\varepsilon''$ etc.
a veris $\alpha,\alpha',\alpha''$ etc. negligere liceat. In hoc itaque casu functio $X''$ non erit
$X$ ipsa, cum qua tamen in eo convenit, quod omnibus $\mu$ valoribus propositis satisfacit.
Omnes autem similes functiones his valoribus satisfacientes sub forma
$X''+Y(\cos\mu x-\cos\mu a)$ contenti erunt, designante $Y$ functionem indefinitam
arcus $x$ eiusdem formae ut $X$, scilicet a sinubus liberam. Functio $X''$ erit
unica ordinis $\mu-1^{\mathrm{ti}}$: quaevis alia similis valoribus datis satisfaciens ad ordinem
altiorem ascendet. Quodsi talis functio $X'''$ aliunde constaret, omnes
similes functiones, per quos valores dati repraesentantur, in hac quoque forma
$X'''+Y(\cos\mu x-\cos\mu a)$ contenti erunt, poteritque $Y$ ita determinari, ut functio
ad ordinem $\mu-1^{\mathrm{tum}}$ depressa prodeat, quae erit ipsa $X''$. Methodus vero facillima
ad hunc finem videtur esse, si primo ex $X'''$ derivetur $X'$ per art. 24 atque
hinc $X''$ per art. 31. Sit itaque terminus quicunque in $X'''$ $L\cos\lambda x$, ponaturque
$\lambda=k\mu+\lambda'$, ita ut sit $\lambda'<\mu$ atque $k$ integer. Tunc per art. 24 pro
illo termino scribi debebit in $X'$
\[
\begin{aligned}
&L\cos k\mu a\cos\lambda'x-L\sin k\mu a\sin\lambda'x,
\\[5pt]
&\qquad\text{si }\lambda'<\tfrac12\mu,\quad\text{neque vero }\lambda'=0\\[5pt]
&L\cos(k+1)\mu a\cos(\mu-\lambda')x+L\sin(k+1)\mu a\sin(\mu-\lambda')x.
\\[5pt]
&\qquad\text{si }\lambda'>\tfrac12\mu\\[5pt]
&L\cos\tfrac12\mu a\cos(k+\tfrac12)\mu a\cos\tfrac12\mu x\\[5pt]
&\qquad{}+L\sin\tfrac12\mu a\cos(k+\tfrac12)\mu a\sin\tfrac12\mu x,
\\[5pt]
&\qquad\text{si }\lambda'=\tfrac12\mu
\end{aligned}
\]
Hinc autem prodit in $X''$, per art. praec. in casu primo
\[
\begin{aligned}
&L\cos k\mu a\cos\lambda'x+L\sin k\mu a\operatorname{cotang}\mu a\cos\lambda'x\\[5pt]
&\qquad{}-L\sin k\mu a\operatorname{cosec}\mu a\cos(\mu-\lambda')x
\end{aligned}
\]
in casu secundo
\[
\begin{aligned}
&L\cos(k+1)\mu a\cos(\mu-\lambda')x-L\sin(k+1)\mu a\operatorname{cotg}\mu a\cos(\mu-\lambda')x\\[5pt]
&\qquad+L\sin(k+1)\mu a\operatorname{cosec}\mu a\cos\lambda'x
\end{aligned}
\]
\SourceLeaf{324}
in casu tertio
\[
\begin{aligned}
&L\cos\tfrac12\mu a\cos(k+\tfrac12)\mu a\cos\tfrac12\mu x\\[5pt]
&\qquad{}-L\sin\tfrac12\mu a\cos(k+\tfrac12)\mu a\operatorname{cotg}\mu a\cos\tfrac12\mu x\\[5pt]
&\qquad+L\sin\tfrac12\mu a\cos(k+\tfrac12)\mu a\operatorname{cosec}\mu a\cos\tfrac12\mu x
\end{aligned}
\]
quae expressio in omnibus tribus casibus reducitur ad
\[
\begin{aligned}
&\frac{L\sin(k+1)\mu a\cos\lambda'x-L\sin k\mu a\cos(\mu-\lambda')x}{\sin\mu a}
\end{aligned}
\]
Quoties autem $\lambda'=0$, habemus in $X'$ simpliciter $L\cos k\mu a$, qui terminus
sine variatione in $X''$ retinetur.

\InterpolationArticle{33}
Sit secundo $X$ functio formae
\[
\begin{aligned}
&\beta'\sin x+\beta''\sin2x+\beta'''\sin3x+\text{etc.}+\beta^n\sin nx
\end{aligned}
\]
unde esse debebit $\mu=n$. Hic fit, ex art. 30, coëfficiens ipsius $A$ in formula
secunda art. 16 pro $T$
\[
\begin{aligned}
&=\frac{1}{\mu\sin\mu a}\times\frac{\sin t}{\sin a}(\sin\mu a+2\sin(\mu-1)a\cos t\\[5pt]
&\qquad{}+2\sin(\mu-2)a\cos2t+\text{etc.}\\[5pt]
&\qquad+2\sin a\cos(\mu-1)t)\\[5pt]
&=\frac{1}{\mu\sin\mu a}(2\cos(\mu-1)a\sin t+2\cos(\mu-2)a\sin2t\\[5pt]
&\qquad{}+2\cos(\mu-3)a\sin3t+\text{etc.}\\[5pt]
&\qquad+2\cos a\sin(\mu-1)t+\sin\mu t)
\end{aligned}
\]
Prorsus similes expressiones pro coëfficientibus ipsorum $B,C,D$ etc. prodeunt,
mutato tantummodo arcu $a$ in $b,c,d$ etc. Quamobrem faciendo
\[
\begin{aligned}
&\zeta'=\frac{2}{\mu\sin\mu a}\{A\cos(\mu-1)a+B\cos(\mu-1)b\\[5pt]
&\qquad{}+C\cos(\mu-1)c+D\cos(\mu-1)d+\text{etc.}\}\\[5pt]
&\zeta''=\frac{2}{\mu\sin\mu a}\{A\cos(\mu-2)a+B\cos(\mu-2)b\\[5pt]
&\qquad{}+C\cos(\mu-2)c+D\cos(\mu-2)d+\text{etc.}\}\\[5pt]
&\zeta'''=\frac{2}{\mu\sin\mu a}\{A\cos(\mu-3)a+B\cos(\mu-3)b\\[5pt]
&\qquad{}+C\cos(\mu-3)c+D\cos(\mu-3)d+\text{etc.}\}\\[5pt]
&\text{etc.}\\[5pt]
&\zeta^{\mu-1}=\frac{2}{\mu\sin\mu a}\{A\cos a+B\cos b+C\cos c+D\cos d+\text{etc.}\}\\[5pt]
&\zeta^\mu=\frac{1}{\mu\sin\mu a}\{A+B+C+D+\text{etc.}\}
\end{aligned}
\]
erit
\[
\begin{aligned}
&T=\zeta'\sin t+\zeta''\sin2t+\zeta'''\sin3t+\text{etc.}+\zeta^\mu\sin\mu t
\end{aligned}
\]
Quum haec formula generaliter pro valore quocunque ipsius $t$ valeat, necessario
\SourceLeaf{325}
cum $X$ identica erit. mutando $t$ in $x$, unde coëfficientes $\zeta',\zeta'',\zeta'''$ etc. ipsis $\beta',\beta'',\beta'''$ etc. resp. aequales erunt.

\InterpolationArticle{34}
Si ut in art. 31 omnibus valoribus propositis per functionem talem
\[
\begin{aligned}
&X'={}\gamma+\gamma'\cos x+\gamma''\cos 2x+\text{etc.}+\gamma^{m}\cos mx\\[5pt]
&+\delta'\sin x+\delta''\sin 2x+\text{etc.}+\delta^{m}\sin mx
\end{aligned}
\]
satisfactum est, existente $m=\tfrac{1}{2}\mu-\tfrac{1}{2}$ vel $=\tfrac{1}{2}\mu$, prout $\mu$ impar vel par est, comparatio huius functionis cum hac
\[
\begin{aligned}
&X''=\zeta'\sin x+\zeta''\sin 2x+\zeta'''\sin 3x+\text{etc.}+\zeta^{\mu}\sin\mu x
\end{aligned}
\]
quam in art. praec. eruimus, hasce aequationes suppeditant:
\[
\begin{aligned}
&\zeta^{\mu}=\gamma\operatorname{cosec}\mu a\\[5pt]
&\zeta'=\delta'+\gamma'\operatorname{cotg}\mu a,\\[5pt]
&\zeta^{\mu-1}=\gamma'\operatorname{cosec}\mu a\\[5pt]
&\zeta''=\delta''+\gamma''\operatorname{cotg}\mu a,\\[5pt]
&\zeta^{\mu-2}=\gamma''\operatorname{cosec}\mu a\\[5pt]
&\zeta'''=\delta'''+\gamma'''\operatorname{cotg}\mu a,\\[5pt]
&\zeta^{\mu-3}=\gamma'''\operatorname{cosec}\mu a\\[5pt]
&\text{etc. usque ad}\\[5pt]
&(\tfrac{1}{2})\zeta^{m}=\delta^{m}+\gamma^{m}\operatorname{cotg}\mu a,\\[5pt]
&(\tfrac{1}{2})\zeta^{\mu-m}=\gamma^{m}\operatorname{cosec}\mu a
\end{aligned}
\]
ubi factor $\tfrac{1}{2}$ coëfficientibus $\zeta^{m},\zeta^{\mu-m}$ praepositus pro eo tantum casu valet, ubi $\mu$ par est, pro altero vero, ubi $\mu$ impar est, omitti debet. In casu itaque posteriore ubi $\mu-m=m+1$, pro quovis coëfficiente $\zeta',\zeta''$ etc. valorem unum per $\gamma,\gamma',\gamma''$ etc. et $\delta',\delta''$ etc. habemus; in priore vero, ubi $\mu-m=m$, pro coëfficiente $\zeta^{\frac{1}{2}\mu}$ duos aequales, pro quibus etiam valor $\delta^{\frac{1}{2}\mu}+\gamma^{\frac{1}{2}\mu}\operatorname{cotg}\mu a+\gamma^{\frac{1}{2}\mu}\operatorname{cosec}\mu a$ adoptari potest. Hinc colligitur
\[
\begin{aligned}
&X''={}X'-\gamma-\gamma'\cos x-\gamma''\cos 2x-\text{etc.}-\gamma^{m}\cos mx\\[5pt]
&+\operatorname{cotg}\mu a(\gamma'\sin x+\gamma''\sin 2x+\gamma'''\sin 3x+\text{etc.}+\gamma^{m}\sin mx)\\[5pt]
&+\operatorname{cosec}\mu a(\gamma\sin\mu x+\gamma'\sin(\mu-1)x\\[5pt]
&\qquad{}+\gamma''\sin(\mu-2)x+\text{etc.}+\gamma^{\mu-m}\sin(\mu-m)x).
\end{aligned}
\]
Quamobrem ex $X'$ producitur $X''$, scribendo pro $\gamma$, $\gamma\operatorname{cosec}\mu a\sin\mu x=\frac{\gamma\sin\mu x}{\sin\mu a}$ et, pro quovis termino $\gamma^{\lambda}\cos\lambda x$
\SourceLeaf{326}
\[
\begin{aligned}
&\gamma^{\lambda}\operatorname{cotg}\mu a\sin\lambda x\\[5pt]
&\qquad{}+\gamma^{\lambda}\operatorname{cosec}\mu a\sin(\mu-\lambda)x
\\[5pt]
&\text{sive}\quad
\frac{\gamma^{\lambda}\cos\mu a\sin\lambda x+\gamma^{\lambda}\sin(\mu-\lambda)x}{\sin\mu a}
\end{aligned}
\]
quod praeceptum lectores cum iis, quae in art. 24 tradidimus, ipsi comparent.

\InterpolationArticle{35}
Si functio $X$ cum $\sin nx$ non abrumpitur, sed ulterius excurrit, terminis sequentibus per $\delta^{\mu+1}\sin(\mu+1)x+\delta^{\mu+2}\sin(\mu+2)x$ etc. expressis, habebimus:
\[
\begin{aligned}
&\zeta'={}\frac{1}{\sin\mu a}\{\delta'\sin\mu a+\delta^{\mu+1}\sin 2\mu a+\delta^{2\mu+1}\sin 3\mu a+\text{etc.}\\[5pt]
&\qquad+\delta^{2\mu-1}\sin\mu a+\delta^{3\mu-1}\sin 2\mu a+\delta^{4\mu-1}\sin 3\mu a+\text{etc.}\}\\[5pt]
&\zeta''={}\frac{1}{\sin\mu a}\{\delta''\sin\mu a+\delta^{\mu+2}\sin 2\mu a+\delta^{2\mu+2}\sin 3\mu a+\text{etc.}\\[5pt]
&\qquad+\delta^{2\mu-2}\sin\mu a+\delta^{3\mu-2}\sin 2\mu a+\delta^{4\mu-2}\sin 3\mu a+\text{etc.}\}\\[5pt]
&\zeta'''={}\frac{1}{\sin\mu a}\{\delta'''\sin\mu a+\delta^{\mu+3}\sin 2\mu a+\delta^{2\mu+3}\sin 3\mu a+\text{etc.}\\[5pt]
&\qquad+\delta^{2\mu-3}\sin\mu a+\delta^{3\mu-3}\sin 2\mu a+\delta^{4\mu-3}\sin 3\mu a+\text{etc.}\}
\end{aligned}
\]
etc. Pro coëfficiente ultimo autem
\[
\begin{aligned}
&\zeta^{\mu}=\frac{1}{\sin\mu a}(\delta^{\mu}\sin\mu a+\delta^{2\mu}\sin 2\mu a+\delta^{3\mu}\sin 3\mu a+\text{etc.})
\end{aligned}
\]
Hae formulae ostendunt, quatenus differentia inter $X$ et $X''$ negligi possit. Haec posterior formula simplicissima erit inter omnes similes, per quas $\mu$ valoribus propositis satisfit; hae vero omnes in formula $X''+Y\sin x(\cos\mu x-\cos\mu a)$ contentae erunt, designante $Y$, ut in art. 32 functionem indefinitam arcus $x$ a sinubus liberam. Et generaliter, si $X'''$ est functio quaecunque eiusdem formae ut $X$, i. e. solos sinus continens, per quam $\mu$ valoribus datis satisfit, formula $X'''+Y\sin x(\cos\mu x-\cos\mu a)$ omnes huiusmodi functiones continebit, quae si $Y$ rite determinatur, ad ordinem $\mu^{\mathrm{tum}}$ deprimi potest, quo pacto necessario functio $X''$ ipsa prodire debet. Prorsus simili modo ut in art. 32 regula generalis sequens ad hunc finem eruitur: Pro quovis termino in $X'''$ tali $L\sin\lambda x$, ubi $\lambda$ est maior quam $\mu$, substituere oportet in $X''$, faciendo $\lambda=k\mu+\lambda'$, ita ut $k\mu$ sit multiplum ipsius $\mu$ proxime minus quam $\lambda$ adeoque $\lambda'$ inter limites $1$ et $\mu$ incl. situs, terminos
\[
\begin{aligned}
&L\frac{\sin(k+1)\mu a}{\sin\mu a}\sin\lambda'x+L\frac{\sin k\mu a}{\sin\mu a}\sin(\mu-\lambda')x
\end{aligned}
\]
qui, quoties fit $\lambda'=\mu$, ad unum $L\frac{\sin\lambda a}{\sin\mu a}\sin\mu x$ reducuntur.
\SourceLeaf{327}
\InterpolationArticle{36}
Transformationes in art. praec. atque in art. 32 traditae concinnius ex theoremate quodam generali deduci possunt, quod quum per se quoque satis elegans sit, paucis hic adhuc attingemus.

\textsc{Theorema.} \textit{Designantibus $\lambda,\lambda',\lambda''$ numeros integros quoscunque. $\mu$ numerum integrum, qui differentias inter illos, $\lambda'-\lambda,\lambda''-\lambda',\lambda-\lambda''$ metitur (e. g. unitatem), $x$ arcum indefinitum, $a$ arcum definitum: functiones}
\[
\begin{aligned}
&P=\sin(\lambda'-\lambda'')a\cos\lambda x+\sin(\lambda''-\lambda)a\cos\lambda'x+\sin(\lambda-\lambda')a\cos\lambda''x\\[5pt]
&Q=\sin(\lambda'-\lambda'')a\sin\lambda x+\sin(\lambda''-\lambda)a\sin\lambda'x+\sin(\lambda-\lambda')a\cos\lambda''x
\end{aligned}
\]
\textit{per} $\cos\mu x-\cos\mu a$ \textit{erunt divisibiles.}

\textit{Demonstr.} Quando $\lambda$ per $\mu$ divisibilis est, ideoque etiam $\lambda',\lambda''$ per $\mu$ divisibiles erunt, facile confirmatur, valorem ipsarum $P,Q$, si substituatur $x=a$, esse identice $=0$; quare $P$ non mutabitur, si pro
\[
\begin{array}{lll}
\cos\lambda x&\text{substituitur}&\cos\lambda x-\cos\lambda a\\
\cos\lambda'x&&\cos\lambda'x-\cos\lambda'a\\
\cos\lambda''x&&\cos\lambda''x-\cos\lambda''a
\end{array}
\]
neque $Q$, si pro
\[
\begin{array}{lll}
\sin\lambda x&\text{substituitur}&\displaystyle\sin\lambda x-\frac{\sin\lambda a\sin\mu x}{\sin\mu a}\\[6pt]
\sin\lambda'x&&\displaystyle\sin\lambda'x-\frac{\sin\lambda'a\sin\mu x}{\sin\mu a}\\[6pt]
\sin\lambda''x&&\displaystyle\sin\lambda''x-\frac{\sin\lambda''a\sin\mu x}{\sin\mu a}
\end{array}
\]
Sed hae sex expressiones per $\cos\mu x-\cos\mu a$ divisibiles sunt, quod pro valore positivo ipsius $\lambda$ de prima et quarta ostendisse sufficit. Scilicet facile per multiplicationem confirmatur, esse
\[
\begin{aligned}
&\sin\mu a(\cos\lambda x-\cos\lambda a)\\[5pt]
&=(\cos\mu x-\cos\mu a)\{2\sin\mu a\cos(\lambda-\mu)x+2\sin 2\mu a\cos(\lambda-2\mu)x\\[5pt]
&\quad \qquad+2\sin 3\mu a\cos(\lambda-3\mu)x+\text{etc.}+2\sin(\lambda-\mu)a\cos\mu x+\sin\lambda a\}
\end{aligned}
\]
\[
\begin{aligned}
&\sin\mu a\sin\lambda x-\sin\lambda a\sin\mu x\\[5pt]
&=(\cos\mu x-\cos\mu a)\{2\sin\mu a\sin(\lambda-\mu)x+2\sin 2\mu a\sin(\lambda-2\mu)x\\[5pt]
&\quad \qquad+2\sin 3\mu a\sin(\lambda-3\mu)x+\text{etc.}+2\sin(\lambda-\mu)a\sin\mu x\}
\end{aligned}
\]
Casus, ubi $\lambda$ est negativus, ad hunc sponte reducitur. Hinc patet, functiones
\SourceLeaf{328}
$P$, $Q$ ex partibus per $\cos\mu x-\cos\mu a$ divisibiles compositas, ideoque ipsas quoque per hunc divisorem divisibiles esse.

II. Quando $\lambda$ per $\mu$ non est divisibilis, sit $l$ numerus integer arbitrarius per $\mu$ divisibilis, ponaturque $\lambda=l+\theta$, $\lambda'=l'+\theta$, $\lambda''=l''+\theta$, unde etiam $l'$, $l''$ per $\mu$ divisibiles erunt. Iam patet, si ponatur
\[
\begin{aligned}
&\sin(l'-l'')a\cos lx+\sin(l''-l)a\cos l'x+\sin(l-l')a\cos l''x=P'\\[5pt]
&\sin(l'-l'')a\sin lx+\sin(l''-l)a\sin l'x+\sin(l-l')a\sin l''x=Q'
\end{aligned}
\]
fieri
\[
\begin{aligned}
&P=P'\cos\theta x-\theta'\sin\theta x,\quad Q=P'\sin\theta x+Q'\cos\theta x,
\end{aligned}
\]
atque functiones $P'$, $Q'$, quippe quae sub casum primum iam absolutum pertinent, per $\cos\mu x-\cos\mu a$ divisibiles: hinc manifesto etiam $P$ et $Q$ per $\cos\mu x-\cos\mu a$ divisibiles erunt. Q.~E.~D. Ceterum demonstratio casus primi ita perfecta est, ut non sine quibusdam explicationibus applicari possit, quoties $\sin\mu a=0$; tunc vero fit $P=0$, $Q=0$, ita ut demonstratione omnino non opus sit.

Quodsi itaque ponitur $\lambda-\lambda'=k\mu$, $\lambda-\lambda''=(k+1)\mu$, patet, per $\cos\mu x-\cos\mu a$ divisibiles esse
\[
\begin{aligned}
&\sin\mu a\cos\lambda x-\sin(k+1)\mu a\cos\lambda'x+\sin k\mu a\cos(\lambda'-\mu)x\\[5pt]
&\sin\mu a\sin\lambda x-\sin(k+1)\mu a\sin\lambda'x+\sin k\mu a\sin(\lambda'-\mu)x
\end{aligned}
\]
adeoque etiam
\[
\begin{aligned}
&\displaystyle\cos\lambda x-\frac{\sin(k+1)\mu a\cos\lambda'x-\sin k\mu a\cos(\mu-\lambda')x}{\sin\mu a}\\[5pt]
&\displaystyle\sin\lambda x-\frac{\sin(k+1)\mu a\sin\lambda'x+\sin k\mu a\sin(\mu-\lambda')x}{\sin\mu a}
\end{aligned}
\]
unde ratio substitutionum in artt. 32, 35 statim elucet: quotiens enim ex divisione posteriore e solis sinubus constabit, adeoque manifesto denuo per $\sin x$ divisibilis erit.

\InterpolationArticle{37}

In artt. 30---36 supposuimus, $\sin\mu a$ non esse $=0$: superest itaque, ut easdem disquisitiones pro eo casu resumamus, ubi $\sin\mu a=0$. Hic statim supponemus, esse $a=0$, vel $a=\frac{180^{\mathord{0}}}{\mu}$.
\SourceLeaf{329}
\par Sit primo $X$ functio formae
\[
\begin{aligned}
&\gamma+\gamma'\cos x+\gamma''\cos2x+\mathrm{etc.}+\gamma^n\cos nx
\end{aligned}
\]
adeoque multitudo coëfficientium incognitorum $=1+n$. Iam scimus, ex omnibus valoribus propositis functionis $X$ eos reiici debere, qui respondent valori ipsius maiori quam $180^{\mathord{0}}$: quatuor itaque casus hic sunt distinguendi:

1) quando $\mu$ par, atque $a=0$, erit $180^{\mathord{0}}$ valor $\frac12\mu+1^{\mathrm{tus}}$ ipsius $x$; quare quum sequentes reiici debeant, remanent valores $\frac12\mu+1$. Hinc esse debebit $n=\frac12\mu$.

2) quando $\mu$ par, atque $a=\frac{180^{\mathord{0}}}{\mu}$. valor $\frac12\mu^{\mathrm{tus}}$ ipsius $x$ erit $180^{\mathord{0}}-\frac{180^{\mathord{0}}}{\mu}$; sequentes, qui fiunt maiores quam $180^{\mathord{0}}$, reiiciendi sunt. Hinc esse debebit $n=\frac12\mu-1$.

3) quando $\mu$ impar est, atque $a=0$, fit valor $\frac12\mu+\frac12^{\mathrm{tus}}=180^{\mathord{0}}-\frac{180^{\mathord{0}}}{\mu}$, et

4) quando $\mu$ impar est, atque $a=\frac{180^{\mathord{0}}}{\mu}$, fit valor $\frac12\mu+\frac12^{\mathrm{tus}}=180^{\mathord{0}}$: sequentes in utroque casu reici debent, adeoque erit $n=\frac12\mu-\frac12$.

Iam quoniam methodus in praecc. adhibita ad casum praesentem, ubi pars valorum datorum a periodo completa antea rescindenda esset, non sine quibusdam ambagibus applicari posset, methodum sequentem praeferimus.

Si per praecepta artt. 20, 22 functio formae
\[
\begin{aligned}
&\gamma+\gamma'\cos x+\gamma''\cos2x+\mathrm{etc.}+\gamma^m\cos mx\\[5pt]
&{}+\delta'\sin x+\delta''\sin2x+\mathrm{etc.}+\delta^m\sin mx
\end{aligned}
\]
investigatur, per quam omnibus $\mu$ valoribus datis satisfit. et in qua $m=\frac12\mu-\frac12$ vel $=\frac12\mu$, prout $\mu$ impar est vel par, coëfficientes $\delta'$, $\delta''$, $\delta'''$ etc. sponte fient $=0$. Nullo enim negotio patet, in expressione tali
\[
\begin{aligned}
&A\sin\lambda a+B\sin\lambda b+C\sin\lambda c+D\sin\lambda d+\mathrm{etc.}
\end{aligned}
\]
fieri vel partem primam $=0$, atque ultimam $=-B\sin\lambda b$, penultimam $=-C\sin\lambda c$, antepenultimam $=-D\sin\lambda d$ etc. puta quando $a=0$; vel ultimam $=-A\sin\lambda a$, penultimam $=-B\sin\lambda b$, antepenultimam $=-C\sin\lambda c$ etc., quando $a=\frac{180^{\mathord{0}}}{\mu}$, quum pro talibus valoribus ipsius $x$, quorum alter alterius complementum ad $360^{\mathord{0}}$ est, valores functionis $X$ aequales sint. Quamobrem functio
\[
\begin{aligned}
&X'=\gamma+\gamma'\cos x+\gamma''\cos2x+\mathrm{etc,}+\gamma^m\cos mx
\end{aligned}
\]
\SourceLeaf{330}
in qua coëfficientes $\gamma$, $\gamma'$, $\gamma''$ etc. determinantur per formulas
\[
\begin{aligned}
&\gamma = \frac{1}{\mu}(A+B+C+D+\mathrm{etc.})\\[5pt]
&\gamma' = \frac{2}{\mu}(A\cos a+B\cos b+C\cos c+D\cos d+\mathrm{etc.})\\[5pt]
&\gamma'' = \frac{2}{\mu}(A\cos 2a+B\cos 2b+C\cos 2c+D\cos 2d+\mathrm{etc.})
\end{aligned}
\]
etc., ultimus autem, quando $\mu$ par est atque adeo $m=\frac12\mu$, per hanc
\[
\begin{aligned}
&\gamma^m=\frac{1}{\mu}(A\cos ma+B\cos mb+C\cos mc+D\cos md+\mathrm{etc.})
\end{aligned}
\]
necessario cum functione $X$ identica erit, siquidem haec non est gradus altioris
quam supra definivimus. Namque in casibus $3$ et $4$ $X$ est ordinis
$\frac12\mu-\frac12{}^{\mathrm{ti}}$, i. e.
eiusdem ut $X'$ et proin per art. 24 cum $X'$ identica. In casu primo $X$ est eiusdem ordinis ut $X'$ et in casu secundo non maioris, quare tum hic tum illic omnes
termini saltem usque ad ordinem $\frac12\mu-1^{\mathrm{tum}}$ in utraque functione convenient
(art. 24). Terminos ordinis $\frac12\mu^{\mathrm{ti}}$ in his functionibus quoque convenire debere.
inde per eundem art.24 patet, quod in $X$ aequatio conditionalis $K\sin ma=L\cos ma$
locum habet; scilicet fit $L=0$, atque in casu primo $\sin ma=0$, in secundo,
ubi $X$ ad ordinem $\frac12\mu-1$ tantummodo ascendit, $K=0$. Ceterum in casu
secundo $X'$ ordinis altioris esse videtur quam $X$, sed in hoc casu terminus ordinis $\frac12\mu^{\mathrm{ti}}$ in $X'$ quoque evanescit, quum fiat
\[
\begin{aligned}
&\gamma^{\frac12\mu}=\frac{1}{\mu}(A\cos 90^{{\mathord{0}}}+B\cos 270^{{\mathord{0}}}+C\cos 450^{{\mathord{0}}}+D\cos 630^{{\mathord{0}}}+\mathrm{etc.})=0
\end{aligned}
\]
ita ut in hoc quoque casu $X'$ revera sit ordinis $m-1^{\mathrm{ti}}$ sive $\frac12\mu-1^{\mathrm{ti}}$.

\InterpolationArticle{38}

Si functio $X$ cum termino $\cos nx$ non abrumpitur, sed ulterius excurrit:
denotatis terminis sequentibus per $\alpha^{n+1}\cos(n+1)x+\alpha^{n+2}\cos(n+2)x+\mathrm{etc.}$ erit
per artt. 21, 23
\[
\begin{aligned}
&\gamma = \alpha\pm\alpha^{\mu}+\alpha^{2\mu}\pm\mathrm{etc.}\\[5pt]
&\gamma' = \alpha'\pm\alpha^{\mu-1}\pm\alpha^{\mu+1}+\alpha^{2\mu-1}+\alpha^{2\mu+1}\pm\mathrm{etc.}\\[5pt]
&\gamma'' = \alpha''\pm\alpha^{\mu-2}\pm\alpha^{\mu+2}+\alpha^{2\mu-2}+\alpha^{2\mu+2}\pm\mathrm{etc.}\\[5pt]
&\gamma''' = \alpha'''\pm\alpha^{\mu-3}\pm\alpha^{\mu+3}+\alpha^{2\mu-3}+\alpha^{2\mu+3}\pm\mathrm{etc.}
\end{aligned}
\]
et sic porro usque ad ultimum $\gamma^m$, quando $\mu$ impar est, vel ad penultimum
$\gamma^{m-1}$, quando $\mu$ par est; signum inferius hic valet, quoties $a=\frac{180^{{\mathord{0}}}}{\mu}$, adeoque
\SourceLeaf{331}
in casu $2$ et $4$. superius in casu $1$ et $2$; denique pro ultimo habebitur in casu
primo
\[
\begin{aligned}
&\gamma^m=\alpha^m+\alpha^{3m}+\alpha^{5m}+\mathrm{etc.}
\end{aligned}
\]
in casu secundo autem $\gamma^m=0$. Hae aequationes ostendunt, quatenus differentiam inter functiones $X$ et $X'$ negligere permissum esse possit. Haec posterior
functio inter omnes, quae $\mu$ valoribus propositis satisfaciunt, simplicissima erit,
quae omnes sub forma $X'+Y\sin(\frac12\mu x-\frac12\mu a)$ contenti erunt, quae in casu
$1$ et $3$ ad $X'+Y\sin\frac12\mu x$, in casibus $2$ et $4$ vero ad $X'+Y\cos\frac12\mu x$ reducitur.
Manifesto autem, si haec expressio eiusdem formae esse debet ut $X$ i. e. a sinubus libera, $Y$ esse debet
\[
\begin{aligned}
&\text{in casu primo formae} \quad  g\sin x+g'\sin 2x+g''\sin 3x+\mathrm{etc.}\\[5pt]
&\text{in casu secundo formae} \quad  g+g'\cos x+g''\cos 2x+\mathrm{etc.}\\[5pt]
&\text{in casu tertio formae} \quad  g\sin\frac12 x+g'\sin\frac32 x+g''\sin\frac52 x+\mathrm{etc.}\\[5pt]
&\text{in casu quarto formae} \quad  g\cos\frac12 x+g'\cos\frac32 x+g''\cos\frac52 x+\mathrm{etc.}
\end{aligned}
\]
Et generalius, designante $X''$ functionem quamcunque ipsi $X$ similem. quae $\mu$ valoribus propositis satisfacit, omnes huiusmodi formae sub formula $X''+Y\sin\frac12\mu x$
vel $X''+Y\cos\frac12\mu x$ contentae erunt, ubi $Y$ functionem indefinitam eius, quam
modo docuimus formae designat. Hoc ita perficere licet, ut sic functio ad ordinem $\frac12\mu$, $\frac12\mu-1$, $\frac12\mu-\frac12$, $\frac12\mu-\frac12$ depressa prodeat, quae manifesto cum $X'$
identica erit. Regula autem generalis pro reductione talis functionis $X''$ ad $X'$
ex art. 24 facile deducitur. Pro quovis termino $L\cos\lambda x$ in $X''$ substitui debet
in $X'$, facto $\lambda=k\mu\pm\lambda'$, ita ut $\lambda'$ non sit maior quam $\frac12\mu$, terminus $\pm L\cos\lambda'x$,
ubi signum inferius accipiendum est, quoties simul $a=\frac{180^{{\mathord{0}}}}{\mu}$ atque $k$ par, superius in casibus reliquis; denique quoties in casu secundo, i. e. pro $a=\frac{180^{{\mathord{0}}}}{\mu}$
et valore pari ipsius $\mu$, evadit $\lambda'=\frac12\mu$, pro $L\cos\lambda x$ statim poni debet $0$ in
$X'$, sive terminus ille omnino negligi.

\InterpolationArticle{39}

Si secundo functio $X$ est formae
\[
\begin{aligned}
&\beta'\sin x+\beta''\sin 2x+\beta'''\sin 3x+\mathrm{etc.}+\beta^n\sin nx
\end{aligned}
\]
adeoque multitudo coëfficientium incognitorum $=n$, etiam multitudo valorum
